\subsection{連立一次方程式}

\subsubsection{行列による表現}

$n$ 個の未知数をもつ $n$ 本の連立一次方程式は
\begin{equation}
 \begin{cases}
  a_{11}x_1+\cdots+a_{1n}x_n=b_1,\\
  \hfill\vdots\\
  a_{n1}x_1+\cdots+a_{nn}x_n=b_n
 \end{cases}
 \quad\Longleftrightarrow\quad
 A\vect{x}=\vect{b}
\end{equation}
と書ける。ここで $A=(a_{ij})$ は\textbf{係数行列}、
$\vect{x}=(x_1,\ldots,x_n)^{\mathsf T}$ は未知数ベクトル、
$\vect{b}=(b_1,\ldots,b_n)^{\mathsf T}$ は定数ベクトルである。
$\det A\ne0$ なら逆行列が存在し、
\begin{equation}
 \boxed{\vect{x}=A^{-1}\vect{b}}
\end{equation}
となる。ただし、実際の数値計算では以下のクラメルの公式や
掃き出し法を使う方が計算の見通しがよいことが多い。

\subsubsection{クラメルの公式}

まず2元連立一次方程式
\begin{equation}
 \begin{cases}
  a_{11}x_1+a_{12}x_2=b_1, &(1)\\
  a_{21}x_1+a_{22}x_2=b_2. &(2)
 \end{cases}
\end{equation}
を加減法で解く。$a_{22}\times(1)-a_{12}\times(2)$ と
$a_{11}\times(2)-a_{21}\times(1)$ を作ると、
\begin{align}
 (a_{11}a_{22}-a_{12}a_{21})x_1&=b_1a_{22}-a_{12}b_2,\label{eq:cramer-x1}\\
 (a_{11}a_{22}-a_{12}a_{21})x_2&=a_{11}b_2-b_1a_{21}.\label{eq:cramer-x2}
\end{align}
これは行列式を用いて
\begin{align}
 \begin{vmatrix}a_{11}&a_{12}\\a_{21}&a_{22}\end{vmatrix}x_1
 &=\begin{vmatrix}b_1&a_{12}\\b_2&a_{22}\end{vmatrix},&
 \begin{vmatrix}a_{11}&a_{12}\\a_{21}&a_{22}\end{vmatrix}x_2
 &=\begin{vmatrix}a_{11}&b_1\\a_{21}&b_2\end{vmatrix}
\end{align}
と表せる。そこで
\[
 D=\begin{vmatrix}a_{11}&a_{12}\\a_{21}&a_{22}\end{vmatrix},\quad
 D_1=\begin{vmatrix}b_1&a_{12}\\b_2&a_{22}\end{vmatrix},\quad
 D_2=\begin{vmatrix}a_{11}&b_1\\a_{21}&b_2\end{vmatrix}
\]
とおけば、$D\ne0$ のとき
\begin{equation}
 \boxed{x_1=\frac{D_1}{D},\qquad x_2=\frac{D_2}{D}}
\end{equation}
を得る。これを\textbf{クラメルの公式}という。

一般に $A\vect{x}=\vect{b}$ に対し、$D=\det A$ とする。
$A$ の第 $k$ 列だけを $\vect{b}$ に置き換えた行列の行列式を
$D_k$ とすれば、$D\ne0$ のとき
\begin{equation}
 \boxed{x_k=\frac{D_k}{D}\quad(k=1,2,\ldots,n)}
\end{equation}
である。未知数が少ないときには便利だが、次数が大きい場合は
多数の行列式が必要になるため、掃き出し法の方が効率的である。

\subsubsection{解の個数と幾何学的意味}

2元一次方程式の各式は $x_1x_2$ 平面上の直線を表す。
したがって、連立方程式を解くことは2直線の共通点を求めることに
等しい。

\begin{figure}[htbp]
\centering
\begin{tikzpicture}[>=Stealth,scale=.9]
 \begin{scope}
  \draw[->] (-.2,0)--(3.7,0) node[right] {$x_1$};
  \draw[->] (0,-.3)--(0,3.2) node[above] {$x_2$};
  \draw[blue,thick] (.2,2.8)--(3.3,.2);
  \draw[red,thick] (.7,-.2)--(2.8,3.1);
  \fill (1.75,1.45) circle (2pt) node[right] {$P$};
  \node at (1.8,-.8) {$D\ne0$：一意解};
 \end{scope}
 \begin{scope}[xshift=5.2cm]
  \draw[->] (-.2,0)--(3.7,0) node[right] {$x_1$};
  \draw[->] (0,-.3)--(0,3.2) node[above] {$x_2$};
  \draw[blue,thick] (.2,.5)--(3.3,2.8);
  \draw[red,thick,dashed] (.2,.5)--(3.3,2.8);
  \node at (1.8,-.8) {$D=0$：無数の解};
 \end{scope}
 \begin{scope}[xshift=10.4cm]
  \draw[->] (-.2,0)--(3.7,0) node[right] {$x_1$};
  \draw[->] (0,-.3)--(0,3.2) node[above] {$x_2$};
  \draw[blue,thick] (.2,.3)--(3.3,2.5);
  \draw[red,thick] (.2,1)--(3.3,3.2);
  \node at (1.8,-.8) {$D=0$：解なし};
 \end{scope}
\end{tikzpicture}
\caption{2元連立一次方程式の三つの解のパターン}
\end{figure}

まとめると次のようになる。
\begin{enumerate}
 \item $D\ne0$：2直線は1点で交わり、解は一意に決まる。
 \item $D=0$ で、定数項まで含めて一方の式が他方の定数倍：
       2直線は一致し、解は無数にある。
 \item $D=0$ だが、定数項まで含めると比例しない：
       2直線は平行で、解はない。
\end{enumerate}
後二者を区別するには係数行列だけでなく定数項も調べる必要がある。
一般の連立方程式では、係数行列 $A$ と拡大係数行列
$[A\mid\vect{b}]$ の階数を用いて
\[
 \operatorname{rank}A=\operatorname{rank}[A\mid\vect{b}]
\]
なら解が存在し、この共通の階数が未知数の個数と等しければ一意解、
小さければ無数の解となる。階数が異なれば解は存在しない。

\subsubsection{掃き出し法（ガウス・ジョルダン法）}

連立方程式の係数と定数項を並べた
\begin{equation}
 [A\mid\vect{b}]
\end{equation}
を\textbf{拡大係数行列}という。次の三つの基本変形は方程式の解を
変えない。
\begin{enumerate}
 \item 二つの行を入れ替える。
 \item 一つの行を $0$ でない定数倍する。
 \item 一つの行に、別の行の定数倍を加える。
\end{enumerate}
これらを使って係数部分を単位行列にする手続きを掃き出し法という。

\begin{example}[掃き出し法による連立方程式の解法]
\[
 \begin{cases}
  3x+2y+z=1,\\
  x+4y-3z=-3,\\
  2x+3y+2z=5
 \end{cases}
\]
を解く。分数を避けて計算を見やすくするため、まず第1行と第2行を
交換する。
\begin{align*}
 \left(\begin{array}{ccc|c}3&2&1&1\\1&4&-3&-3\\2&3&2&5\end{array}\right)
 &\xrightarrow{R_1\leftrightarrow R_2}
 \left(\begin{array}{ccc|c}1&4&-3&-3\\3&2&1&1\\2&3&2&5\end{array}\right)\\
 &\xrightarrow[,R_3\leftarrow R_3-2R_1,]{R_2\leftarrow R_2-3R_1}
 \left(\begin{array}{ccc|c}1&4&-3&-3\\0&-10&10&10\\0&-5&8&11\end{array}\right)\\
 &\xrightarrow{R_2\leftarrow-\frac1{10}R_2}
 \left(\begin{array}{ccc|c}1&4&-3&-3\\0&1&-1&-1\\0&-5&8&11\end{array}\right)\\
 &\xrightarrow[,R_1\leftarrow R_1-4R_2,]{R_3\leftarrow R_3+5R_2}
 \left(\begin{array}{ccc|c}1&0&1&1\\0&1&-1&-1\\0&0&3&6\end{array}\right)\\
 &\longrightarrow
 \left(\begin{array}{ccc|c}1&0&0&-1\\0&1&0&1\\0&0&1&2\end{array}\right).
\end{align*}
したがって
\[
 \boxed{x=-1,\qquad y=1,\qquad z=2}.
\]
最後に元の3式へ代入すると、計算間違いを検算できる。
\end{example}

\begin{example}[演習]
次の連立方程式を掃き出し法で解け。ただし、途中で必要なら行を
入れ替えてよい。
\[
 \begin{pmatrix}1&-2&1\\-2&4&1\\3&-5&2\end{pmatrix}
 \begin{pmatrix}x\\y\\z\end{pmatrix}
 =\begin{pmatrix}0\\1\\1\end{pmatrix}.
\]
答えは次のとおりである。
\[
 \boxed{x=\frac73,\qquad y=\frac43,\qquad z=\frac13}
\]
\end{example}

\subsubsection{掃き出し法による逆行列}

$AA^{-1}=I$ である。$A^{-1}$ の各列を未知ベクトルとみれば、これは
$A\vect{x}=\unit{1},\ldots,A\vect{x}=\unit{n}$ という $n$ 組の
連立方程式である。右辺をまとめた拡大行列に同じ行基本変形を行い、
\begin{equation}
 \left(A\mid I\right)\longrightarrow\left(I\mid A^{-1}\right)
\end{equation}
とすれば逆行列を一度に求められる。途中で左側を $I$ にできなければ、
$A$ は正則でなく逆行列は存在しない。

\begin{example}[逆行列の計算]
\[
 A=\begin{pmatrix}2&1&-2\\1&1&-1\\2&3&-1\end{pmatrix}
\]
に対し、$(A\mid I)$ を掃き出す。初めに第1行と第2行を交換すると
計算しやすい。
\begin{align*}
\left(\begin{array}{ccc|ccc}2&1&-2&1&0&0\\1&1&-1&0&1&0\\2&3&-1&0&0&1\end{array}\right)
&\xrightarrow{R_1\leftrightarrow R_2}
\left(\begin{array}{ccc|ccc}1&1&-1&0&1&0\\2&1&-2&1&0&0\\2&3&-1&0&0&1\end{array}\right)\\
&\longrightarrow
\left(\begin{array}{ccc|ccc}1&0&0&2&-5&1\\0&1&0&-1&2&0\\0&0&1&1&-4&1\end{array}\right).
\end{align*}
よって
\begin{equation}
 \boxed{A^{-1}=\begin{pmatrix}2&-5&1\\-1&2&0\\1&-4&1\end{pmatrix}}.
\end{equation}
求めた後に $AA^{-1}=I$ を計算すれば、各成分をまとめて検算できる。
\end{example}
